\documentclass[12pt]{article}
\usepackage[utf8]{inputenc}
\usepackage[T1]{fontenc}
\usepackage{hyperref}
\title{Punan Tubu Language Resources for K-12 Curriculum Development}
\author{Compiled by Claude Sonnet 5.5 (claude-sonnet-5-5) for Andr\'as Kornai}
\date{}
\begin{document}
\maketitle

\section{Introduction}
Punan Tubu is an Austronesian language of Borneo, spoken in the Malinau, Mentarang and Sekatak river drainages of North Kalimantan, Indonesia. One review reports roughly forty settlements and a population of about 3,500 to 4,000; its position within the North Sarawak group is described as uncertain. Searches found academic and archive material only, and no school resources.

\section{Primary Educational Resources}
No primer, textbook or school programme was found. The language is not known to be taught in schools.

\section{Linguistic and Reference Materials}
\begin{itemize}
\item Wikipedia, ``Punan Tubu language''. \url{https://en.wikipedia.org/wiki/Punan_Tubu_language}
\item C\'esard et al. (2017), Wacana 18(1): article/review discussing Punan Tubu, including a reference to a dictionary by a Punan Tubu speaker. \url{https://halshs.archives-ouvertes.fr/halshs-02542193v1/file/2017-C%C3%A9sard%20et%20al-Wacana%2018-1.pdf}
\item Soriente and Inagaki, handout on Kalimantan languages (mentions Punan Tubu documentation work). \url{https://lingdy.aa-ken.jp/wp-content/uploads/2012/01/120217_soriente_inagaki_a.pdf}
\item University of Hawai\textquoteleft i ScholarSpace record of handwritten notes on Punan Tubu dialects. \url{https://scholarspace.manoa.hawaii.edu/items/39366f64-8869-4bb2-b0a5-12153576a925/full}
\end{itemize}

\section{Religious and Community Resources}
\begin{itemize}
\item Joshua Project, Punan Tubu in Indonesia: people-group profile. \url{https://joshuaproject.net/people_groups/print/14481/ID}
\end{itemize}

\section{Digital and Audio Resources}
\begin{itemize}
\item MPI-EVA Jakarta Field Station archive page listing Punan Tubu data sessions (access reported as restricted). \url{https://lingweb.eva.mpg.de/archive/jakarta/data_PUNDUL.php.html}
\end{itemize}

\section{Curriculum Development Context}
Materials are scholarly and not classroom-ready. Any curriculum would depend on contacting the researchers and the community, and on the restricted archival recordings, with Indonesian as the likely school language.
\end{document}